% =============================================================================
\section{Binding Code}
\label{sec:bind-code}
% =============================================================================

% -----------------------------------------------------------------------------
\subsection{Return value policy}
\label{sec:bind-code:rvp}
% -----------------------------------------------------------------------------
Return value policy, kernel specific:

In some cases, in order to assure proper behaviour and memory
management of a return value from a function, explicitly specifying a
return value policy is needed. The policy may depend on whether the
returned object is a newly created object, a reference to some already
existing (internal) object or a const reference to some existing
object. [Additionally in some cases we need to tie the lifetime of the
result to the lifetime of the arguments, or the other way around]

For some kernel functions the kind of return value depends on the
kernel (epec vs epic), and as such the return value policy needs to be
changed accordingly. This is why for these functions a
\myLstinline{Kernel_return_value_policy} type is being passed as the
return value policy. The type is defined differently (via preprocessor
macros) based on the kernel being exposed.

% -----------------------------------------------------------------------------
\subsection{Scopes}
\label{sec:bind-code:scopes}
% -----------------------------------------------------------------------------
All functions and classes related to some package were exposed inside
a \myLstinline{export_[package]()} function. All such calls are
happening in \myLstinline{export_module.cpp}. We wrap each call with a
local (c++) scope and call the \myLstinline{SET_SCOPE(module_name)}
macro, which causes every exposed function/class inside this local
scope to be exposed inside the \myLstinline{module_name} submodule of
the CGALPY module. This prevents name clashes of classes with the same
name in different packages.

% -----------------------------------------------------------------------------
\subsection{Using Templates to avoid code duplication}
\label{sec:bind-code:templates}
% -----------------------------------------------------------------------------
Some types have the exact same methods exposed (for example
\myLstinline{Triangulation_2} and
\myLstinline{Delaunay_triangulation_2}. In those cases we use a
templated function (i.e \myLstinline{export_triangulation}) templated
with the type to expose and call the binding functions there.

% -----------------------------------------------------------------------------
\subsection{SFINAE}
\label{sec:bind-code:sfinae}
% -----------------------------------------------------------------------------
\url{https://en.cppreference.com/w/cpp/language/sfinae}
\url{https://doc.cgal.org/latest/Kernel_23/group__do__intersect__linear__grp.html}

\myLstinline{CGAL::do_intersect} has an overloaded implementation for
many supported types, however, not every combination of two supported
types is implemented.

For conveniently binding all overloaded versions of
\myLstinline{CGAL::do_intersect}, we used SFINAE: This is how a
version overloaded with a combination of \myLstinline{T1} and
\myLstinline{T2} is bound: First \myLstinline{bind_do_intersect()}
calls \myLstinline{bind_do_intersect_1T} templated with some supported
type {T1}. \myLstinline{bind_do_intersect_1T()}, templated
with \myLstinline{T1}, calls \myLstinline{bind_do_intersect_2T}
templated with \myLstinline{T1} and some supported type
\myLstinline{T2}, and a \myLstinline{bool} argument.

The declaration of \myLstinline{bind_do_intersect_2T}, templated with
both \myLstinline{T1} and \myLstinline{T2} is only valid when
templated with a combination of types for which an overloaded
implementation of \myLstinline{CGAL::do_intersect} exists. Otherwise,
due to SFINAE, the blank implementation template of
\myLstinline{bind_do_intersect_2T} expecting (...) as an argument
will be invoked (it would normally not be invoked as variadic
parameters have the lowest rank for the purpose of overload
resolution).

\url{https://en.cppreference.com/w/cpp/language/variadic_arguments}

% -----------------------------------------------------------------------------
\subsection{Arrangement Extension}
\label{sec:bind-code:ae}
% -----------------------------------------------------------------------------
The arrangement is always extended with a generic python object. This
allows for flexibility and helps dealing with not being able to use a
custom data type as a template parameter.

% -----------------------------------------------------------------------------
\subsection{Arrangement Overlay traits}
\label{sec:bind-code:overlay}
% -----------------------------------------------------------------------------
Since a custom c++ \myLstinline{overlay_traits} class cannot be
defined from python, we instead expose a new class called
\myLstinline{Arr_python_overlay_traits}. This class is a model of the
OverlayTraits concept.
\url{https://doc.cgal.org/latest/Arrangement_on_surface_2/classOverlayTraits.html}

It accepts in its constructor 10 python functions, each accepting two
python objects (the extended arrangement data type) and returning a
python object. Those 10 functions are saved as class members and are
then called in the respective 10 [overloaded] versions of
\myLstinline{create_vertex}, \myLstinline{create_edge},
\myLstinline{create_face object}) (the matching is done by
the order of the passed functions and the order of the [overloaded]
versions in the official concept reference).

This way the user can still define the way the result face/edge/vertex
data will be calculated.

% -----------------------------------------------------------------------------
\subsection{kd tree}
\label{sec:bind-code:kdtree}
% -----------------------------------------------------------------------------
The k-d tree dimension and the dimension of related classes is
templated and can thus only be decided at compile time via the cmake
variable \myLstinline{CGALPY_SPATIAL_SEARCHING_DIMENSION}.

In order to assert the dimension the bindings were compiled with, a
free function \myLstinline{get_spatial_searching_dimension()} was
exposed that returns the dimension as an \myLstinline{int}.

% -----------------------------------------------------------------------------
\subsection{\myLstinline{Python_distance}}
\label{sec:bind-code:distance}
% -----------------------------------------------------------------------------
Similarly to overlay traits, a user may want to define a custom
distance class to use with search queries. Here too, as defining a new
c++ class is not possible from python, we provide a similar solution:
A \myLstinline{Python_distance} class is exposed, which is a model of
\myLstinline{GeneralDistance} concept.
\url{https://doc.cgal.org/5.2/Spatial_searching/classGeneralDistance.html}
It expects 5 functions in its constructor, and those are saved as
class members and are used for the respective 5 operations [the
optional are not included] of the concept.

Making some things more pythonic:

% -----------------------------------------------------------------------------
\subsection{Iterators and circulators}
\label{sec:bind-code:iters}
% -----------------------------------------------------------------------------
For binding iterators we used the \myLstinline{bind_iterator}
templated function and its variations, in order to allow pythonic
iteration. Circulators too are artificially converted to iterators
supporting \myLstinline{__next__} and throwing a stop iteration when
reaching the first item again, by being wrapped in a templated
"\myLstinline{Iterator_from_circulator}" object which is the one being
exposed to python.

% -----------------------------------------------------------------------------
\subsection{Functions populating an output iterator}
\label{sec:bind-code:fnc-out}
% -----------------------------------------------------------------------------
Instead, a wrapper that expects a python list and populates it with
the results is being exposed.

% -----------------------------------------------------------------------------
\subsection{Functions expecting begin and end iterators}
\label{sec:bind-code:fnc-range}
% -----------------------------------------------------------------------------
Instead, a wrapper that expects a python list as an argument and calls
the function with the list's begin and end iterators is exposed.
